\documentclass[10pt,a4paper]{article}

\usepackage[utf8]{inputenc}
\usepackage[T1]{fontenc}
\usepackage{lmodern}
\usepackage[a4paper,left=17mm,right=17mm,top=15mm,bottom=15mm]{geometry}
\usepackage{enumitem}
\usepackage{titlesec}
\usepackage[hidelinks,unicode]{hyperref}
\input{glyphtounicode}
\pdfgentounicode=1

\hypersetup{
  pdfauthor={Bruno Henrique Nunes},
  pdftitle={Bruno Henrique Nunes - Curriculo de IA Aplicada e Machine Learning},
  pdfsubject={IA aplicada, aprendizado de maquina, visao computacional e ciencia de dados}
}

\pagestyle{empty}
\setlength{\parindent}{0pt}
\setlength{\parskip}{2pt}
\setlength{\emergencystretch}{2em}
\hyphenpenalty=10000
\exhyphenpenalty=10000
\setcounter{secnumdepth}{0}
\titleformat{\section}{\normalsize\bfseries}{}{0pt}{}
\titlespacing*{\section}{0pt}{7pt}{3pt}
\setlist[itemize]{leftmargin=1.2em,itemsep=1.5pt,topsep=2pt,parsep=0pt,partopsep=0pt}
\newcommand{\role}[3]{\textbf{#1 --- #2}\hfill {\small #3}\par}

\begin{document}
\fontsize{10}{12.2}\selectfont

{\Large\bfseries Bruno Henrique Nunes}\par
{\large\bfseries Engenheiro de IA Aplicada e Machine Learning | Ciência de Dados}\par
Brasil | Aberto a oportunidades remotas que contratam no Brasil\par
\href{mailto:bruno.nunes.1987@gmail.com}{bruno.nunes.1987@gmail.com}
\quad\textbar\quad \href{https://github.com/bhnunes}{github.com/bhnunes}
\quad\textbar\quad \href{https://www.linkedin.com/in/brunohenriquenunes/}{linkedin.com/in/brunohenriquenunes}\par

\section{Resumo Profissional}
Engenheiro de IA aplicada e Machine Learning, com mestrado concluído em Inteligência Artificial pela Universidade Estadual de Campinas (UNICAMP). Desenvolvo agentes de IA, fluxos multiagentes e aplicações com LLMs, RAG, NLP e engenharia de prompts. Combino desenvolvimento de APIs em Python, avaliação de LLMs, testes automatizados e observabilidade com pesquisa reprodutível em deep learning e análise estatística.

\section{Competências Técnicas}
\textbf{IA aplicada e engenharia de LLMs:} Python, agentes de IA, sistemas multiagentes, Large Language Models (LLMs), Retrieval-Augmented Generation (RAG), Processamento de Linguagem Natural (NLP), engenharia de prompts (prompt engineering), avaliação de LLMs, busca semântica e embeddings de texto.\par
\textbf{Machine Learning e Data Science:} PyTorch, deep learning, visão computacional, análise estatística, desenho de experimentos, intervalos de confiança, qualidade de dados, modelagem de tópicos e Apache Spark.\par
\textbf{Engenharia de software:} FastAPI, APIs REST, testes automatizados, observabilidade, AWS Lambda, AWS Simple Email Service (SES), Git e Azure DevOps CI/CD.

\section{Experiência Profissional}
\role{Cientista de Dados (Data Scientist)}{IBM}{Ago 2026 -- Atual}
\begin{itemize}
  \item Desenvolvo agentes de IA em Python e fluxos multiagentes para casos de uso de IA aplicada, engenharia de software e Data Science.
  \item Desenvolvo aplicações com LLMs fundamentadas em recuperação, usando RAG, NLP, recuperação de evidências, prompt engineering e APIs REST com Python/FastAPI.
  \item Implemento defesas contra prompt injection, guardrails de entrada e saída, testes automatizados e observabilidade; apoio a implantação de aplicações.
\end{itemize}

\role{Professor Assistente}{UNICAMP}{Jul 2024 -- Dez 2025}
Dei suporte a mais de 200 estudantes de graduação em gestão de bancos de dados e programação orientada a objetos; conduzi sessões de resolução de problemas em SQL, álgebra relacional e OOP e apoiei avaliações. Vínculo de estágio.

\role{Desenvolvedor de Software Sênior}{PAteam}{Jan 2024 -- Mai 2024}
Desenvolvi automações de ingestão de dados em Python usando AWS Lambda, APIs REST e AWS SES.\par
\vspace{2pt}
\role{Desenvolvedor de Software Sênior}{Expleo Group}{Mai 2022 -- Ago 2023}
\begin{itemize}
  \item Desenvolvi automações em Python/Robocorp e integrações de APIs para migração de sistemas legados para Salesforce, usando APIs REST do Salesforce e SQL Server; documentei fluxos e apoiei CI/CD no Azure DevOps, monitoramento e resolução de problemas.
\end{itemize}
\vspace{1pt}
\role{Desenvolvedor de Software Sênior}{Nestl\'e}{Fev 2022 -- Ago 2023}
\begin{itemize}
  \item Entreguei 13 automações corporativas com Python, Blue Prism, PowerShell, SQL e Automation Anywhere; revisei código em mais de 10 projetos, orientei desenvolvedores e apoiei a aprovação de mudanças para mais de 50 robôs.
\end{itemize}
\vspace{1pt}
\role{Desenvolvedor de Software Sênior}{Deloitte}{Nov 2021 -- Fev 2022}
Transformei requisitos de stakeholders em automações com Python, Node.js e UiPath para um projeto de integração de IoT na área da saúde; apresentei soluções a clientes e fornecedores.\par
\vspace{2pt}
\textbf{IBM --- Cargos anteriores}\par
Desenvolvedor de RPA (out. 2019 -- nov. 2021); Desenvolvedor de Software (out. 2018 -- out. 2019); Analista de TI (jan. 2017 -- out. 2018); Especialista de Suporte (jan. 2012 -- out. 2014).\par
\begin{itemize}
  \item Entreguei automações e integrações com Python/Selenium, APIs, SAP, Db2, Cloudant, Blue Prism, VBA e JavaScript.
  \item Criei análises de faturas com VBA e Power Query para mais de 20.000 registros e 10 indicadores de validação; o tempo de processamento reportado caiu de mais de três horas para cerca de seis minutos.
  \item Usei SQL em verificações de conformidade e relatórios de ajustes de cobrança; prestei suporte em inglês a clientes dos Estados Unidos e do Canadá.
\end{itemize}

\newpage
\section{Pesquisa e Projetos Selecionados}
\textbf{Deep Learning para Histopatologia da Próstata}\hfill {\small Fev 2024 -- Set 2026}\par
Pesquisa de mestrado, Universidade Estadual de Campinas (UNICAMP)\par
\begin{itemize}
  \item Desenvolvi um pipeline de pesquisa para segmentação binária de câncer de próstata em imagens histopatológicas, com modelagem baseada em CNNs e Transformers e PyTorch.
  \item Implementei preparação e controle de qualidade dos dados, particionamento por paciente, treinamento de modelos, otimização de ensemble e avaliação final.
  \item Preservei configurações, proveniência dos dados e resultados experimentais para apoiar a reprodutibilidade. O depósito na REDU contém artefatos de pesquisa gerados; as imagens originais do DIAGSET não são redistribuídas.
\end{itemize}
{\small \href{https://github.com/bhnunes/Master_Coding}{Código: github.com/bhnunes/Master\_Coding}\quad\textbar\quad
\href{https://doi.org/10.25824/redu/EDPCF2}{Artefatos: doi.org/10.25824/redu/EDPCF2}}\par
\vspace{2pt}
\textbf{GEODITE --- Potenciais Interatômicos Equivariantes}\par
Contribuí para a implementação, o treinamento e a avaliação de um modelo de passagem de mensagens equivariante para potenciais interatômicos. Coautor do preprint \emph{Equivariant Interatomic Potentials without Tensor Products} (arXiv, janeiro de 2026).\par
{\small \href{https://arxiv.org/abs/2601.15492}{arxiv.org/abs/2601.15492}}\par

\section{Projetos Selecionados de IA Aplicada e Ciência de Dados (2026)}
\begin{itemize}
  \item \textbf{Qualidade de respostas de LLMs:} Combinei métricas lexicais, similaridade semântica, NLP e embeddings de texto com avaliação assistida por modelos em casos ambíguos.
  \item \textbf{Engenharia e avaliação de agentes de IA:} Desenvolvi ferramentas em Python para fluxos multiagentes de engenharia de software, com orientações sensíveis ao repositório e testes automatizados; projetei análises estatísticas pareadas com teste de Wilcoxon, intervalos de confiança bootstrap e medidas de uso de tokens.
  \item \textbf{Data Science aplicada e NLP:} Prototipei análises apoiadas por RAG, classificação multilíngue de texto e busca semântica; apliquei validação de dados, modelagem de tópicos, desenho de métricas estatísticas, Apache Spark e enriquecimento de telemetria com APIs REST, PostgreSQL e Kubernetes.
\end{itemize}

\section{Formação Acadêmica}
\textbf{Mestrado em Inteligência Artificial --- Universidade Estadual de Campinas (UNICAMP)}\hfill {\small 2024 -- 2026}\par
Concluído em setembro de 2026; pesquisa em deep learning para histopatologia da próstata.\par
\textbf{MBA em Engenharia de Software --- FIA Business School}\hfill {\small 2022 -- 2023}\par
\textbf{Bacharelado em Engenharia Física / Física Aplicada --- Universidade de São Paulo}\hfill {\small 2006 -- 2010}\par

\section{Certificações Selecionadas}
AWS Certified Generative AI Developer -- Professional (out. 2026)\par
AWS Certified Machine Learning Engineer -- Associate (jul. 2026)\par
AWS Certified Developer -- Associate (out. 2025)\par
AWS Certified Solutions Architect -- Associate (abr. 2025)

\section{Idiomas}
\textbf{Português:} nativo.\quad\textbf{Inglês:} proficiência profissional plena; TOEIC Listening and Reading: 890 (set. 2025).\par
\textbf{Espanhol e italiano:} básico.

\end{document}
